\section{Conclusion}
We present VEXMLM, a vocabulary-extended variant of XLM-R designed to address excessive subword fragmentation in Ge'ez-script languages through: (i) language-specific SentencePiece tokenization for Amharic and Tigrinya; (ii) augmentation with \vocabAdded{} Ge'ez-derived subword tokens; (iii) mean-based embedding initialization for compatibility with the pretrained embedding space; and (iv) a two-stage adaptation framework in which the vocabulary-extended model is first adapted through continued MLM pretraining on Amharic and Tigrinya corpora and then fine-tuned for downstream tasks.

VEXMLM reduces tokenizer fertility relative to XLM-R by \tokFertRedAmh\% for Amharic and \tokFertRedTir\% for Tigrinya, below that of Glot500. Downstream, it improves named entity recognition over XLM-R, is comparable on sentiment analysis, and scores below XLM-R on extractive question answering. The ablation indicates that continued MLM pretraining is the primary driver of the gains: vocabulary expansion alone lowers accuracy on out-of-vocabulary words, and mean-based initialization is only a slightly better starting point than random initialization.

Overall, our findings show that language-aware vocabulary expansion makes the tokenization of underrepresented Ge'ez-script languages substantially more efficient, but that its downstream benefit is task-dependent and requires adapting the new embeddings through continued pretraining.
